\documentclass[11pt]{article}

\usepackage{sectsty}
\usepackage{graphicx}
\usepackage{amsmath,amssymb}
% Margins
\topmargin=-0.45in
\evensidemargin=0in
\oddsidemargin=0in
\textwidth=6.5in
\textheight=9.0in
\headsep=0.25in

\title{A Review of \\ One-step targeted maximum likelihood estimation for
	targeting cause-specific absolute risks and survival curves}
%\author{}
%\date{\today}

\begin{document}
	\maketitle	
	% Optional TOC
	% \tableofcontents
	% \pagebreak
	
This article presents an extension of two authors' independent works based on the targeted maximum likelihood estimation proposed by the second author and his coauthor \cite{1} and the work of both authors and their co-author \cite{2} on a one-step targeted maximum likelihood estimation methodology in a single parameter survival modeling on cause-specific hazards to a multi-parameter scenario. 


Unlike in the single-parameter case in which the authors used an iterative targeting procedure that updates cause-specific hazards recursively to solve the efficient influence curve equation for that particular univariate parameter, in this article, the authors defined a universal least favorable sub-model for the cause-specific hazards and a corresponding one-step targeting procedure that allows us to solve multiple score equations simultaneously. 

The article highlights two major developments:  (1) a simultaneous inference for treatment effects on multiple absolute risks in competing risk settings, and (2) an inference for the full survival curve or absolute risk curve across time. 

Asymptotic results related to the new methodology were presented. The proposed methodology can also be applied to settings with right-censored event times and can handle competing risk scenarios.

%%%%%%
\begin{thebibliography}{1}
	%
	\bibitem{1}
	 van der Laan M. J., Rubin D. B. (2006). Targeted maximum likelihood
	 learning \emph{Int J Biostat.} \textbf{2(1)}, article 11.
	% 
	
	\bibitem{2}
	 Rytgaard, H. C.W., Eriksson, F., van der Laan, M. J. (2023). Estimation of time-specific intervention effects on continuously distributed time-to-event outcomes by targeted maximum likelihood estimation. \emph{Biometrics} \textit{https://doi.org/10.1111/biom.13856}.
	%
\end{thebibliography}
	
	
	
	
\end{document}





